\section{Discussion}
\label{sec:discussion}

\paragraph{Crop-level splitting inflates accuracy, and the inflation is leakage.} With one
registered test, sized by simulation, Study~2 detects what Study~1's six separate tests could
not: on SIPaKMeD, crop-level splitting raises accuracy by \HoneEst~pp, consistently across
twelve architectures from ResNet-18 to DINOv2. The factorial removes the obvious alternative.
At equal crop and field counts, giving the model the test fields' siblings raises accuracy by
\HthreeEst~pp, and the effect replicated on fresh seeds (\HthreeRepEst~pp), while giving it 43\%
more distinct fields without leakage does nothing. The natural gap is therefore leakage, not
the extra training diversity that crop-level splits also bring.

\paragraph{The size of the inflation follows the structure.} SIPaKMeD releases about four crops
per field and every field belongs to one class; CRIC releases about 29 cells per slide image,
and slide identity carries most of the class information (Section~\ref{sec:structure}). The
inflation scales accordingly, from \HoneEst~pp (an error rate \SipErrorRatio\ times higher at
the parent level) to \HtwoEst~pp (\CricErrorRatio\ times), with balanced accuracy on CRIC
falling by \CricBaccEst~pp. The closed forms of Section~\ref{sec:structure} predict which
benchmarks are exposed before any model is trained.

\paragraph{How the leak is used.} The two registered mechanism tests of Study~2 and the eight
follow-up tests converge on a description by elimination. The leaked advantage does not need a
background: it survives masking every background in training and testing (F3) and removing it
at test time, and the background-transplant route was small and uncertain (H4). It is not
memorisation of shared pixels: it is as large for test cells that share no pixel with any
training crop (exploratory). It is not a slide-level label prior: on CRIC, seen slides helped
cells whose class disagreed with the slide's majority as much as cells that agreed, and pulled
them \emph{away} from the majority class (F1). And it is not carried by a separable,
linearly decodable code for the field: erasing that code from the features in closed form,
or suppressing it with an adversary that cannot hide it from a linear probe, cut
recoverability by \RecoveryCutMin--\RecoveryCutMax\% of its baseline without moving the gap (F5, F7, F8). What remains is that the
advantage depends on the cell being rendered as its own field renders cells: transferring a
test cell's colour statistics to another field of the same class removed about three quarters
of it (F2). The simplest account consistent with all of this is that a model trained on a
field's siblings calibrates its class evidence to that field's appearance (stain, colour
balance, and the texture they produce in the cell itself), so cells from a seen field sit
further inside their class region. Random colour jitter during training did not prevent this
(F4), presumably because jitter perturbs each image around its own appearance and leaves the
field-specific structure of that appearance learnable.

\paragraph{Decodable is not used, and saliency is not use either.} Study~1 established
\emph{availability}: a probe recognised fields at 109 times chance. Study~2 shows that
availability and use come apart. The field identity a linear probe can read can be removed
almost entirely without touching the inflation, in line with evidence from language models that
probe-decodable information need not be what a model relies on~\citep{elazar2021amnesic,ravichander2021probing}.
Grad-CAM, the tool most often used to argue that a medical classifier looks at the right
thing, showed leak and clean models attending to the same places; a saliency audit would have
cleared both. Neither probes nor saliency maps can certify that an evaluation is free of
hierarchical leakage; only the split can.

\paragraph{What survives, and what does not fix it.} Within a benchmark, crop-level splitting
preserved the architecture ranking reasonably well (Kendall $\tau=\KendallTau$), so relative
comparisons under one protocol are not meaningless. Absolute numbers are another matter: results
reported for the same dataset under crop-level and parent-level splits are not comparable, and on
CRIC the protocol difference is several times larger than the spread between the six
architectures under either protocol. Neither of the usual robustness measures removed the
inflation: background removal left it intact (F3), strong colour augmentation left it intact
(F4), and stain normalisation reduced recoverability only slightly while costing accuracy (F6).

\subsection{Limitations}
\label{sec:limitations}

\paragraph{Parents are not patients.} SIPaKMeD fields, Cx22 source images, and CRIC slide
images are the deepest units the releases identify. Several fields can come from one
patient, so parent-disjoint evaluation is necessary but may not be sufficient for
clinical independence, and every estimate here is a lower bound on what patient-level
grouping would remove.

\paragraph{The factorial trains on fewer crops.} Holding crop and field counts equal
across leak and clean cells requires 220 crops per class (1,100 in total, against 3,239
in Phase~1), so the 2C and 2B effects are measured on smaller training sets than the
headline contrast. The decomposition that transfers them to the natural splits assumes
the leakage effect scales with the leak rate and the diversity effect scales with the
logarithm of the field-count ratio; it is labelled approximate.

\paragraph{Transplants are imperfect interventions.} The foreground mask is a fixed
threshold rule, inpainting leaves a smoothed background, and feathering mixes a few
pixels of donor and recipient. The \emph{self} condition measures these artifacts, and
every contrast compares conditions that share them, but composites are still off the
training manifold.

\paragraph{Gradient reversal removes what its head can see.} The adversary suppresses
parent identity that a two-layer head can read from pooled features; information that
remains decodable by another probe, or that enters through spatial features before
pooling, can survive. Recoverability is therefore measured with an independent probe,
and H5 is declared untestable rather than refuted if suppression does not work.

\paragraph{One recipe.} Every model uses the same fixed fine-tuning recipe. Stronger
augmentation, longer training, or stain normalisation could change both the size of
the inflation and the models' reliance on acquisition context.

\paragraph{One host.} All runs used one four-GPU server with a four-core CPU allowance,
which is why augmentation runs on the GPU. Wall-clock times are not portable, and the
GPU implementation of antialiased resizing differs from PIL's by uint8 rounding.

\paragraph{No clinical claim.} The study audits research benchmarks. It does not
estimate screening sensitivity, specificity, prevalence-calibrated performance, or
effects on readers.

\subsection{Recommendations for dataset releases and evaluations}
\label{sec:recommendations}

\begin{enumerate}
\item \textbf{Release the hierarchy.} Every observation should carry a stable
  observation identifier, its acquisition parent (field, source image, or slide), a
  patient or specimen identifier where ethics and law permit, and, if duplicates were
  audited, a duplicate-component identifier.
\item \textbf{Split at the deepest released unit and publish the split.} Evaluate on
  parent-disjoint (or patient-disjoint) folds and release the exact fold manifest with
  its hash, so that results reported under one dataset name are comparable.
\item \textbf{Report both units.} State observation and parent counts for every split,
  and the fraction of test observations with a training sibling.
\item \textbf{Check the structure before training.} The closed forms of
  Section~\ref{sec:structure} (\texttt{cytofair risk-report}) give the expected straddle,
  sibling rate, diversity ratio, and class--parent dependence of a crop-level split from
  the parent table alone.
\item \textbf{Treat one question as one test.} When the same question is asked of many
  architectures, test it once with a mixed model that treats architectures and splits as
  random effects, and size the study by simulation; report per-architecture contrasts
  descriptively.
\item \textbf{Do not certify a split with a probe, a saliency map, or an augmentation.} Leakage
  inflated accuracy with field identity erased from the features, with identical Grad-CAM maps,
  with backgrounds removed, and with strong colour augmentation. Only a parent-disjoint split
  removes it.
\item \textbf{Keep the evidence.} Save per-sample predictions, final weights, and the
  exact split for every run, so that confusion matrices, calibration, and later
  mechanistic analyses do not require retraining.
\end{enumerate}

